% Display and inline formulas that came out half text, half picture (visual hunt, family dispmath).
% One frame per mechanism; tests/test_classify.py pins what each becomes.
\documentclass{beamer}
\usepackage{amsmath,amssymb}
\DeclareMathOperator{\SSIM}{SSIM}
\newtheorem{claim}{Claim}

\begin{document}

% 1: an inline sum with limits after italic theorem words: its box widens the gap to the
% formula's words past what build_lines joins
\begin{frame}{Inline sum with limits}
\begin{claim}
In a graph, $\sum\limits_{v \in V} d(v) = 2|E|$ holds for all finite graphs.
\end{claim}
\end{frame}

% 2: an inline stacked fraction whose bar is the width of its numerator
\begin{frame}{Inline fraction}
The mean of the first six numbers is $\frac{1+2+3+4+5+6}{6}$ which is three and a half.
\end{frame}

% 3: a display with big delimiters and an integral
\begin{frame}{Big delimiters}
The norm is
\[ \|f\|_p = \left( \int_X |f|^p \, d\mu \right)^{1/p}. \]
for every measurable function.
\end{frame}

% 4: a display integral and a radical in the second line, under words of the first
\begin{frame}{Operators in a paragraph}
We integrate the density of the measure over the whole space and obtain the total mass
$\displaystyle\int_X \rho \, d\mu$ of the body, which is finite.

The spectral gap of a random regular graph is close to the Ramanujan bound given by
$2\sqrt{d-1}$ for large graphs, as Friedman proved.

Hashing into $\sqrt{n}$ buckets balances the two costs.
\end{frame}

% 5: a display with operator names, a minimum and a fraction
\begin{frame}{Operator names}
The structural similarity is
\[ \SSIM(x, y) = \min\left(1, \frac{(2\mu_x\mu_y + c_1)(2\sigma_{xy} + c_2)}{(\mu_x^2 + \mu_y^2 + c_1)(\sigma_x^2 + \sigma_y^2 + c_2)}\right) \]
\end{frame}

% 6: aligned rows that start with a relation, and cases with a condition column
\begin{frame}{Aligned rows and cases}
\begin{align*}
(a + b)^2 &= (a + b)(a + b) \\
&= a^2 + ab + ba + b^2 \\
&= a^2 + 2ab + b^2
\end{align*}
\[ f(x) = \begin{cases} x^2 & \text{if } x \ge 0, \\ -x & \text{otherwise.} \end{cases} \]
\end{frame}

% 7: a display sum whose limits reach below its box
\begin{frame}{Display sum}
\[ \sum_{k=1}^{n} k = \frac{n(n+1)}{2} \]
\end{frame}

% 8: a display integral with nothing under it: its ink hangs an em below every glyph box
\begin{frame}{Display integral}
The mass is
\[ m = \int \rho \, dV \]
\end{frame}

\end{document}
